\section{Дофамин по вспышке света}
\label{sec:dishdopamine}

Единственный известный нам прямой опыт, где культуре дали дофамин как учителя, поставлен на платформе из органоидов, к которой исследователи подключаются удалённо~\cite{Jordan2024}. В жидкость, омывающую органоиды, добавили дофамин в светочувствительной «клетке»: связанная молекула не действует на рецепторы. Концентрация клеточного дофамина --- $1$~мМ. Когда нужно наградить сеть, на неё светят ультрафиолетом: клетка распадается, и дофамин выходит в объём.

Петля устроена так. Один и тот же стимул подаётся раз за разом; если он вызвал больше импульсов, чем прежде, включается вспышка и высвобождается дофамин. За $240$ повторов число вызванных импульсов в целом росло. Авторы сами пишут, что по одному такому опыту нельзя заключить, что рост вызван дофамином~\cite{Jordan2024}.

Для инженера это первая попытка собрать в чашке правило трёх множителей~\eqref{eq:threefactor}: активность отправителя и получателя оставляет след, а общий сигнал решает, закрепить ли изменение. Но сигнал здесь бывает только положительным: награда есть или её нет, отрицательной ошибки $\delta<0$ установка не выдаёт. Дофамин к тому же расплывается и откачивается медленно, как концентрация в~\eqref{eq:lowpass}, так что частые вспышки могут просто поднять его общий уровень. Чтобы отличить обучение от этого, нужны два контроля: столько же вспышек в случайные моменты, не связанные с ответом сети, и те же вспышки без клеточного дофамина. И нужна проверка после отдыха: осталось ли изменение, когда дофамин ушёл.

\section{Дофамин учит кремний}
\label{sec:biohybrid}

В обратном опыте живые клетки учат электронику~\cite{Keene2020}. Клетки, выделяющие дофамин, выращивают в микроканале над органическим нейроморфным элементом --- транзистором, проводимость канала которого играет роль веса синапса. Дофамин, вышедший из клеток, окисляется у электрода элемента, и это меняет проводимость. Устройство повторяет и механизм уборки медиатора из щели, поэтому вес можно не только изменить надолго, но и вернуть обратно.

Схемотехнически это интегратор с утечкой по химическому входу: вес копит поток дофамина, а «откачка» задаёт, как быстро он забывается, --- та же RC-цепь, что и~\eqref{eq:lowpass}. Главное здесь направление связи. Щуп из главы~\ref{sec:debugging} не видит широковещательных регистров, потому что они химические. Этот элемент читает именно химический регистр и превращает его в электрический вес. Для интерфейсов мозг---компьютер это путь к датчику, который слышит не только шину данных, но и регистры управления.

\section{Органоиды со своими дофаминовыми нейронами}
\label{sec:midbrainorg}

Из стволовых клеток человека умеют выращивать органоиды, похожие на ту область мозга, где сидят \nt{DOPA}-нейроны. В них есть работающие дофаминовые нейроны, которые выделяют дофамин~\cite{Jo2016}. Такие органоиды выращивают в основном для изучения болезней. Опытов, где их дофамин служил бы учителем для другой сети, мы не нашли.

Для машины из живых нейронов это недостающая деталь. Стенд из главы~\ref{sec:livingmachine} --- шина данных и управление потоком без регистров управления; здесь источник широковещательного сигнала уже есть. Не хватает критика: сети, которая вычисляла бы ошибку предсказания~\eqref{eq:ctd} и управляла бы этими нейронами. Соединить органоид коры с таким органоидом так, чтобы дофамин выделялся по ошибке, --- открытая задача, и пока это гипотеза, а не опыт.

\section{Органоид балансирует шестом без дофамина}
\label{sec:cartpole}

Самый полный из недавних опытов обходится без дофамина совсем~\cite{Robbins2026}. Органоиды коры мыши подключили в замкнутую петлю к задаче «шест на тележке»: активность органоида двигает тележку, положение шеста подаётся обратно стимуляцией. Между попытками органоид получает короткие высокочастотные учебные стимулы. Какие именно --- выбирает программа обучения с подкреплением.

Результаты измерены:
\begin{itemize}
\item у большинства органоидов учебные сигналы, выбранные программой, дали лучший результат, чем случайно выбранные сигналы или отсутствие сигнала;
\item после $45$ минут отдыха улучшение не сохранялось;
\item блокада рецепторов глутамата обоих типов, AMPA и NMDA, отменяла улучшение.
\end{itemize}

Последний пункт говорит, где живёт выученное: в \nt{GLUT}-синапсах, и через тот рецептор, который в разделе~\ref{sec:weightstore} работает как детектор совпадения входа и выхода. Второй пункт показывает, чего не хватает: быстрое изменение есть, закрепления нет, как у быстрого ученика из главы~\ref{sec:motorlearning} без медленного. А роль учителя изменилась: инженера, который выбирает рисунок тока, заменила программа, но учитель по-прежнему стоит снаружи чашки.

Среди более ранних опытов обучения культур на электродных матрицах мы также не нашли ни одного, где использовался бы дофамин: учебным сигналом везде служила электрическая стимуляция.

\section{Кто учит культуру}

\begin{table}[t]
\centering
\footnotesize
\setlength{\tabcolsep}{4pt}
\renewcommand{\arraystretch}{1.15}
\begin{tabular}{>{\raggedright}p{2.7cm}>{\raggedright}p{2.5cm}>{\raggedright}p{3.2cm}>{\raggedright\arraybackslash}p{4.4cm}}
\toprule
Опыт & Кто учит & Учебный сигнал & Что известно \\
\midrule
остановка стимуляции при нужном ответе & инженер & конец стимуляции & ответ закрепляется --- измерено \\
DishBrain (гл.~\ref{sec:dishbrain}) & инженер & предсказуемый или случайный ток & значимый рост розыгрышей за сеанс только при полной обратной связи, при тишине эффект меньше --- измерено; от сеанса к сеансу неустойчив; причина обсуждается \\
шест на тележке & программа обучения с подкреплением & выбранные программой высокочастотные стимулы & лучше случайных, но не сохраняется после отдыха --- измерено \\
дофамин по вспышке & правило «больше импульсов --- награда» & дофамин, освобождённый светом & рост импульсов --- наблюдение; роль дофамина не доказана \\
биогибридный синапс & живые клетки & дофамин от клеток & долгое изменение и возврат веса элемента --- измерено \\
органоиды с \nt{DOPA}-нейронами & --- & --- & источник дофамина есть; как учитель не использован \\
\bottomrule
\end{tabular}
\caption{Кто и чем учит живые нейроны на чипе.}
\label{tab:dishteachers}
\end{table}

Во всех опытах, где учится сама культура, учитель пока снаружи (таблица~\ref{tab:dishteachers}): инженер, программа или правило, заданное инженером. Дофамин вошёл в чашку лишь однажды, и его роль ещё предстоит доказать. Собрать в чашке настоящий широковещательный канал --- с критиком, ошибкой и следом, как в главе~\ref{sec:dofa}, --- следующий шаг, которого никто ещё не сделал.
